\chapter{Introdução}
\label{cap:introducao}

Circuitos digitais formam a base de praticamente todos os sistemas computacionais modernos, e seu estudo é parte obrigatória da formação em Ciência da Computação. No currículo do curso de Ciência da Computação da \gls{ufsc}, esse conteúdo é introduzido na disciplina de Circuitos e Técnicas Digitais, na qual os estudantes têm o primeiro contato formal com portas lógicas, álgebra booleana, circuitos combinacionais e sequenciais, e com uma \gls{hdl}, tipicamente \gls{vhdl}.

Paralelamente, é comum que esses mesmos estudantes estejam cursando, no mesmo período ou em período próximo, a disciplina de \gls{poo}, que utiliza Python como linguagem principal. O resultado é que o aluno transita, dentro de um curto intervalo de tempo, entre dois paradigmas de programação extremamente distintos: de um lado, Python, uma linguagem imperativa, sequencial, de sintaxe enxuta e com um ciclo de desenvolvimento quase instantâneo (editar, executar, ver o resultado); de outro, VHDL, uma linguagem declarativa e concorrente, com sintaxe verbosa e formal, que descreve \textit{hardware} e não um programa a ser executado passo a passo.

A motivação deste trabalho nasce justamente dessa distância entre as duas experiências de aprendizado vividas pelo estudante de Ciência da Computação. Observam-se, na prática, algumas dificuldades recorrentes:

\begin{itemize}
    \item \textbf{Distância de paradigma e de sintaxe}: enquanto em Python o aluno já está habituado a um código curto e direto, VHDL exige a declaração explícita de entidades, arquiteturas, sinais, tipos e processos antes mesmo de descrever a lógica desejada. Essa sintaxe extensa e formal, combinada com um modelo de execução concorrente pouco intuitivo para quem só programou de forma sequencial, torna os primeiros contatos com a disciplina de Circuitos e Técnicas Digitais particularmente custosos;
    \item \textbf{Ambientes de desenvolvimento pesados e pouco ágeis}: as \glspl{ide} tradicionalmente utilizadas para VHDL (como Quartus, da Intel, ou Vivado, da Xilinx/AMD) são ferramentas robustas, voltadas para síntese em \glspl{fpga} reais, porém pesadas, lentas para abrir e configurar, e com uma curva de aprendizado própria que é somada à da linguagem em si;
    \item \textbf{Tempo de compilação/síntese elevado}: diferentemente do ciclo quase instantâneo de ``salvar e rodar'' ao qual o aluno está acostumado em Python, a simulação ou síntese de um circuito VHDL pode levar segundos a minutos, mesmo para exemplos simples. Esse tempo de espera atrapalha a experimentação rápida e a construção de intuição sobre o comportamento do circuito, que é justamente o tipo de aprendizado que se beneficia de iteração rápida e \textit{feedback} imediato;
    \item \textbf{O problema não é exclusivo do VHDL}: linguagens de descrição de \textit{hardware} aparentadas, como Verilog e SystemVerilog, compartilham a mesma natureza concorrente, sintaxe formal e ecossistema de ferramentas pesado, de modo que a dificuldade relatada aqui não é uma particularidade do VHDL, mas sim das HDLs tradicionais como classe de linguagem.
\end{itemize}

A busca por ferramentas que tornem os conteúdos de \textit{hardware} mais acessíveis a estudantes de Computação não é recente. Há mais de duas décadas, \citeonline{hackersitte2004} propuseram um ambiente interativo para o ensino de lógica digital combinacional e sequencial, motivados pela falta de \textit{software} de apoio adequado a esse ensino. Mais recentemente, \citeonline{garay2019} recorreram à visualização de simulações em VHDL para ensinar um tópico de Arquitetura e Organização de Computadores que, na experiência dos autores, é bastante complexo para os estudantes, evidenciando que mesmo o VHDL tem sido usado como meio, e não como fim, em iniciativas que buscam reduzir barreiras de aprendizado.

Na comunidade brasileira de educação em computação, dois mapeamentos sistemáticos recentes sintetizam parte dessa produção. \citeonline{santini2023} identificaram treze jogos educacionais voltados ao ensino de Circuitos Digitais publicados nos últimos dez anos, partindo da constatação de que ``a disciplina de Circuitos Lógicos muitas vezes é considerada tediosa pelos alunos'' \cite[p.~814]{santini2023}; em relato de experiência anterior, \citeonline{carvalhonakamura2013} já haviam observado falta de interesse dos alunos pela disciplina, vista por eles como um assunto próprio da Engenharia. Já \citeonline[p.~240]{dondici2025} abrem seu mapeamento afirmando que ``o ensino de Circuitos Lógicos, embora essencial nos cursos de Computação, enfrenta desafios pela complexidade dos conceitos''. A complexidade também aparece em tópicos mais específicos da área: \citeonline[p.~191]{jurgina2021} observam que a complexidade do tema confiabilidade de circuitos ``é uma barreira para que seu estudo seja incorporado em sala de aula'', o que motivou o desenvolvimento de uma ferramenta dedicada a esse ensino. Na literatura internacional, \citeonline{haase2022} propõe uma ferramenta de simulação interativa de circuitos digitais, executada no navegador, para apoiar o modelo de sala de aula invertida e o estudo autônomo. A dificuldade também se reflete em dados de evasão: em um estudo sobre os cursos de Computação da Universidade Estadual de Maringá, \citeonline{fukao2023} estimaram taxas de evasão acumuladas de 47,6\% em Ciência da Computação e de 54,5\% em Informática para os ingressantes de 2005 a 2021, e os estudantes evadidos apontaram a didática adotada pelos professores como o principal fator percebido como impactante para a evasão. Em conjunto, essas evidências situam a dificuldade aqui relatada como um problema reconhecido e ativamente pesquisado há pelo menos duas décadas, tanto na literatura internacional quanto na comunidade brasileira de educação em computação.

Diante desse cenário, este trabalho parte da hipótese de que um simulador de circuitos digitais implementado em Python, linguagem com a qual o estudante já possui familiaridade prévia por meio da disciplina de POO, pode servir como uma etapa intermediária de aprendizado, permitindo que os conceitos fundamentais de portas lógicas, propagação de sinais e componentes sequenciais (como registradores) sejam compreendidos em um ambiente leve, de execução imediata e sintaxe já conhecida, antes que o aluno enfrente a barreira adicional imposta pelas HDLs.

\section{Problema}

Como apoiar o aprendizado dos conceitos básicos de circuitos digitais (portas lógicas, propagação de sinais e componentes sequenciais) de forma que o estudante de Ciência da Computação possa se concentrar na lógica do circuito em si, sem que a sintaxe, o ambiente de desenvolvimento e o tempo de compilação de uma HDL tradicional se tornem obstáculos adicionais nesse primeiro contato?

\section{Objetivos}

O objetivo geral deste trabalho é desenvolver, em Python, um simulador didático de circuitos digitais baseado em propagação reativa de sinais, capaz de apoiar o aprendizado dos conceitos introdutórios da disciplina de Circuitos e Técnicas Digitais.

Os objetivos específicos são:

\begin{itemize}
    \item Modelar uma abstração de barramento de \textit{bits} (sinal), capaz de notificar automaticamente os componentes conectados a ele quando seu valor muda;
    \item Implementar as portas lógicas básicas (NOT, AND, OR, XOR) como componentes que recalculam sua saída automaticamente a partir da propagação de mudanças em suas entradas;
    \item Implementar um componente de registrador (\textit{latch}/\textit{flip-flop}) sensível a borda de subida de um sinal de \textit{clock}, permitindo a construção de circuitos sequenciais e de memória;
    \item Demonstrar, a partir desses componentes básicos, a construção de circuitos compostos (por exemplo, um meio-somador, um contador binário e um \textit{flip-flop} do tipo \textit{toggle});
    \item Avaliar, por meio desses circuitos de exemplo, se a abordagem proposta reproduz corretamente o comportamento esperado dos circuitos digitais correspondentes.
\end{itemize}

\section{Justificativa}

Este trabalho se justifica por atacar um obstáculo concreto observado na trajetória de estudantes de Ciência da Computação: o salto abrupto entre a programação em Python (aprendida em POO) e a descrição de \textit{hardware} em VHDL (exigida em Circuitos e Técnicas Digitais). Ao propor uma ferramenta que utiliza uma linguagem já familiar ao aluno para modelar os mesmos conceitos (portas lógicas, propagação de sinais, registradores), pretende-se reduzir a carga cognitiva do primeiro contato com a disciplina, isolando a dificuldade dos conceitos de circuitos digitais da dificuldade adicional imposta pela sintaxe e pelas ferramentas de uma HDL.

Além do valor pedagógico, o desenvolvimento do simulador também tem valor técnico: sua implementação exige a aplicação de conceitos de engenharia de \textit{software} (como o padrão de projeto \textit{Observer} e a programação orientada a eventos) para resolver um problema de domínio distinto (simulação de \textit{hardware}), evidenciando a transferência de conceitos entre áreas do curso de Ciência da Computação.

\section{Estrutura do Trabalho}

O restante deste trabalho está organizado da seguinte forma: o Capítulo~\ref{cap:fundamentacao} apresenta a fundamentação teórica necessária, cobrindo circuitos digitais, HDLs e o modelo de programação reativa/padrão \textit{Observer}; o Capítulo~\ref{cap:trabalhos} discute trabalhos relacionados, incluindo outros simuladores de circuitos digitais e iniciativas pedagógicas para HDLs; o Capítulo~\ref{cap:proposta} descreve a arquitetura e a implementação do simulador; o Capítulo~\ref{cap:resultados} apresenta os circuitos construídos como estudo de caso e discute os resultados obtidos; e o Capítulo~\ref{cap:conclusao} traz as considerações finais e sugestões de trabalhos futuros.
